\section{Du RAG v1 au v1.1}\label{sec:v11}

\subsection{Diagnostic de v1}
Avant de construire l'agent, nous avons corrigé le système existant : sans cela, un gain de l'agent pourrait n'être
que la correction d'un défaut de la référence. Le diagnostic, détaillé dans \texttt{docs/baseline\_v1.1\_report.md}, a
relevé : une réécriture de la requête par le LLM qui inventait des numéros d'articles et ajoutait des mots parasites ;
des frontières de \term{chunks} placées sur des renvois internes ; un \term{reranker} qui échouait sans le signaler ; et une
table de synonymes dont les entrées reprenaient les thèmes des questions d'évaluation, c'est-à-dire une fuite.

\subsection{Les étapes et leur effet}
Le tableau~\ref{tab:steps} donne l'effet de chaque étape sur le jeu de développement. L'étape 3 \emph{baisse} le score :
la table de synonymes donnait une partie des réponses, et sa suppression rend le chiffre honnête. Les étapes 6 et 7,
qui amélioraient une métrique au détriment d'une autre, ont été rejetées.

\begin{table}[H]
\centering
\small
\begin{tabular}{p{7.2cm}rrrrc}
\toprule
Étape (dev, 35 questions) & Hit@6 & Hit@1 & Rappel & MRR & Gardée \\
\midrule
v1 (étiquettes finales) & 57{,}1 & 34{,}3 & 25{,}4 & 0{,}403 & -- \\
1. Découpage aux vrais intitulés, identifiants stables & 71{,}4 & 25{,}7 & 32{,}0 & 0{,}403 & oui \\
2. Question brute d'abord (traduction AR/EN seulement) & 88{,}6 & 54{,}3 & 53{,}5 & 0{,}657 & oui \\
3. Suppression de la table de synonymes (fuite) & 85{,}7 & 45{,}7 & 51{,}3 & 0{,}593 & oui \\
4. Reranker : correction, réserve de 20, int8 & 85{,}7 & 57{,}1 & 52{,}7 & 0{,}651 & oui \\
5. Fusion par rang réciproque (RRF) & 88{,}6 & 57{,}1 & 54{,}1 & 0{,}659 & oui \\
6. Texte de contexte avant l'\emph{embedding} & 82{,}9 & 62{,}9 & 52{,}3 & 0{,}689 & non \\
7. bge-m3 au lieu de Solon & 82{,}9 & 54{,}3 & 51{,}1 & 0{,}643 & non \\
\bottomrule
\end{tabular}
\caption{Les sept étapes de v1 à v1.1, chacune cumulative, mesurées sur le dev. Scores lus dans \texttt{results/} (fichiers \texttt{chunkfix\_dev\_final\_labels} et \texttt{item2} à \texttt{item7}).}
\label{tab:steps}
\end{table}


\subsection{Résultat sur le jeu de test}
Le jeu de test n'a été utilisé qu'après la fin des réglages (tableau~\ref{tab:v1v11}, figure~\ref{fig:v1v11}).
Sur les \NVoneTestN{} questions de test, Hit@6 passe de \NVoneTestHitsix~\% (v1) à \NVoneoneTestHitsix~\% (v1.1) et le MRR
de \NVoneTestMrr{} à \NVoneoneTestMrr. Ces chiffres de v1.1 utilisent l'indice de code fourni par le jeu de données
(variante A1) ; sans cet indice, la référence équitable A2 obtient \NPvrTestATwoHitsix~\% (section~\ref{sec:recup}).

\begin{table}[H]
\centering
\small
\begin{tabular}{lrrrrr}
\toprule
Système & Hit@6 & Hit@1 & Rappel & MRR & n \\
\midrule
v1 (test) & 57{,}7 & 32{,}7 & 34{,}4 & 0{,}405 & 52 \\
v1.1, A1 (test) & 80{,}8 & 59{,}6 & 50{,}4 & 0{,}687 & 52 \\
v1.1, A1 (dev) & 88{,}6 & 57{,}1 & 54{,}1 & 0{,}659 & 35 \\
\bottomrule
\end{tabular}
\caption{Récupération, v1 et v1.1 (A1 : avec l'indice de code du jeu de données). Source : \texttt{results/}.}
\label{tab:v1v11}
\end{table}

\figfile{f2_v1_vs_v11}{v1 et v1.1 sur le jeu de test (\NQTest{} questions), quatre métriques de récupération.}{fig:v1v11}

\subsection{Vitesse}
Pour v1.1, le temps moyen par question sur le test est de \NVoneoneSecTotal{} s (traduction comprise ;
\NVoneoneTestLlmcalls{} appels LLM pour les \NQTest{} questions, uniquement pour traduire l'arabe et l'anglais). Pour v1, la
récupération et le \term{reranking} prennent \NVoneSecRetrieval{} s en moyenne ; le temps LLM enregistré pour v1
(\NVoneSecLlm{} s) n'est pas une mesure fiable, car les appels étaient servis par un cache disque avec des attentes de quota.
La comparaison de vitesse entre v1 et v1.1 ne porte donc que sur la récupération, plus les appels LLM que v1 effectue
avant de chercher (réécriture et sous-requêtes).
